\cotefile{156-9}
\scanpage{74}
\mspage{70}
Ceci posé, \deleted{\unsure{les} $\beta_{\unsure{\ast|\times}}$} \unsure{on aura|on a}
\[
\beta_{\ast}(x')=C\bigl(\underbrace{\bar{\beta}(x')}_{\substack{\uparrow\\ \mathbb{\Sigma}_{E}}}\bigr)=\operatorname{cosupp}_{E}^{\circ}\bigl(\bar{\beta}(x')\bigr)
\]
et l'application $\beta$ \unsure{de|en} \unsure{(ii)|(vi)} satisfait la condition \unsure{(xx)|(ix)} \unsure{ssi} $\beta_{\ast}$ est \unsure{associée|associé} à $\beta$, ou encore
\[
\operatorname{supp}_{E}^{\circ}\beta(x')=C\bigl(\bar{\beta}(x')\bigr)
\]
ou aussi
\[
\operatorname{cosupp}^{\circ}\unsure{C|(}\beta(x')=\bar{\beta}(x').
\]

\uline{Définition} Les $\alpha_{\ast}$ satisfaisant les conditions équivalentes \unsure{(i)} \unsure{(bis)|(ii)} \unsure{(bis$'$)|(ii$'$)} sont appelées correspondances disjonctives \deleted{\unsure{entre}} $\underline{E},\underline{E}'$. La correspondance $\beta_{\ast}$ \unsure{de|du} (i) \inserted{(qui est aussi une corr. disj.)} \unsure{est} appelée \deleted{\unsure{la}} corr. disjonctive \deleted{\uline{\unsure{associée}}} \uline{transposée} de $\alpha_{\ast}$.

\uline{Scholie} On trouve \deleted{ainsi} une bijection
\[
\alpha_{\ast}\mathrel{\unsure{\vdash|\mapsto}}{}^{t}\alpha_{\ast}:\operatorname{Corrdis}(\underline{E},\underline{E}')\longrightarrow\operatorname{Corrdis}(\underline{E}',\underline{E})
\]
et on a
\[
\left\{\begin{aligned}
{}^{t}({}^{t}\alpha_{\ast})&=\alpha_{\ast}\\
{}^{t}({}^{t}\beta_{\ast})&=\beta_{\ast}
\end{aligned}\right.
\]
% archive sheet number: 73
